\documentclass{article}

% Use NeurIPS 2024 style; [final] removes line numbers (required per project guidelines).
% Download from: https://media.neurips.cc/Conferences/NeurIPS2024/Styles.zip
\usepackage[final]{neurips_2024}

\usepackage[utf8]{inputenc}
\usepackage[T1]{fontenc}
\usepackage{hyperref}
\usepackage{url}
\usepackage{booktabs}
\usepackage{amsfonts}
\usepackage{nicefrac}
\usepackage{microtype}
\usepackage{graphicx}
\usepackage{amsmath}
\usepackage{amssymb}

\title{Predicting Keystrokes from Electromyography: Comparing CNN, CNN+RNN, and Transformer Architectures}

\author{
  [Author names and UCLA emails] \\
  Department of Electrical and Computer Engineering \\
  University of California, Los Angeles \\
  \texttt{\{emails\}@ucla.edu} \\
}

\begin{document}

\maketitle

\begin{abstract}
We address the task of decoding surface electromyography (sEMG) signals from the wrists to predict QWERTY keystrokes, using the emg2qwerty dataset and the single-subject setup (ID 89335547). We compare three families of architectures: (1) a time--depth separable (TDS) convolutional baseline, (2) a CNN+RNN hybrid with LSTM or GRU (including a bidirectional variant), and (3) a Transformer encoder, all trained with connectionist temporal classification (CTC) loss and evaluated with character error rate (CER). We describe our data preprocessing, augmentation, and training setup, report validation and test CER for each architecture, and discuss why certain designs yield better or worse performance. Our experiments show that [INSERT: e.g., the CNN+RNN / Transformer outperforms the TDS baseline; or summarize main finding]. We also [INSERT: any additional experiments, e.g., data fraction, channels, or preprocessing ablations if performed].
\end{abstract}

% -----------------------------------------------------------------------------
\section{Introduction}
\label{sec:intro}

Minimizing character error rate (CER) for a single user with limited data is practically important for personalized assistive interfaces: it directly affects typing accuracy and usability. We therefore focus on (1) reducing test CER on the provided single subject (ID 89335547) by comparing multiple architectures beyond the default TDS CNN, and (2) understanding which inductive biases---recurrent vs.\ attention-based temporal modeling---best suit this sequential, unsegmented decoding task.

We use the official train/validation/test splits for subject 89335547 and evaluate architectures that satisfy the project requirement of at least one recurrent model while also exploring a Transformer encoder. Our goal is to determine whether adding recurrence (LSTM/GRU) or self-attention (Transformer) on top of a shared spectrogram-based front-end improves decoding over the TDS baseline, and to report and interpret the results clearly.

% -----------------------------------------------------------------------------
\section{Methods}
\label{sec:methods}

\subsection{Task and Data}
The input is a sequence of sEMG signals from two wrist bands (2 bands $\times$ 16 channels, 2\,kHz). The model predicts a sequence of characters (including a null symbol for ``no key'') over the same time span. Labels are keylogger keystrokes with arbitrary timestamps, so the number of characters does not match the number of time steps; we use CTC loss~\cite{ctc} and evaluate with CER $(S + D + I)/N$ (substitutions, deletions, insertions over total characters $N$).

We use the emg2qwerty dataset~\cite{emg2qwerty} and the pre-defined train/val/test splits for subject 89335547. Each session is stored as HDF5; we use windowed segments for training and validation and full sessions at test time for evaluation.

\subsection{Preprocessing and Augmentation}
Raw sEMG is converted to a log-magnitude spectrogram (e.g., \texttt{LogSpectrogram} with \texttt{n\_fft=64}, \texttt{hop\_length=16}, yielding a 125\,Hz frame rate). Training augmentations include: band rotation (electrode permutation per band), temporal jitter of the alignment between EMG and labels, and SpecAugment (time and frequency masking). Validation and test use the same spectrogram without augmentation. Optionally we use channel-wise or global normalization (e.g., z-score) depending on the chosen transform config.

\subsection{Model Architectures}
All models share a common front-end: per-channel spectrogram normalization (\texttt{SpectrogramNorm}), then a rotation-invariant MLP applied per band (\texttt{MultiBandRotationInvariantMLP}), then flattening to a single feature vector per time step. The temporal encoder and head differ by architecture.

\textbf{TDS CNN (baseline).} A time--depth separable convolutional encoder~\cite{emg2qwerty} (TDS blocks + fully connected blocks) processes the sequence and preserves the feature dimension. A linear layer maps the encoder output to the character vocabulary (including null) followed by log-softmax for CTC.

\textbf{CNN+RNN.} The same CNN front-end (TDS conv blocks only, no FC blocks) is followed by a recurrent layer (LSTM or GRU) with configurable hidden size, number of layers, dropout, and bidirectionality. The RNN output (hidden size, or twice that if bidirectional) is projected to the character set and log-softmax for CTC.

\textbf{Transformer.} The same front-end is followed by a linear projection to a model dimension $d$, sinusoidal positional encoding (computed for the actual sequence length to support long test sessions), and a standard Transformer encoder (multi-head self-attention and feed-forward layers). The encoder output is projected to the character set and log-softmax for CTC.

Hyperparameters (e.g., window length 8000 samples, padding, learning rate schedule) are set via Hydra configs (\texttt{config/model/tds\_conv\_ctc.yaml}, \texttt{conv\_rnn\_ctc.yaml}, \texttt{conv\_bilstm\_ctc.yaml}, \texttt{transformer\_ctc.yaml}).

\subsection{Training}
We use PyTorch Lightning with the Adam optimizer and a learning rate schedule (e.g., linear warmup + cosine annealing). Training is monitored by validation CER; the best checkpoint is selected and used for final validation and test reporting. CTC loss is applied as implemented in the codebase (blank index = null class).

% -----------------------------------------------------------------------------
\section{Results}
\label{sec:results}

Table~\ref{tab:cer} reports validation and test CER for the compared architectures on subject 89335547. [\textbf{Replace the placeholder values with your actual numbers.}]

\begin{table}[ht]
\centering
\caption{Character error rate (\%) on subject 89335547.}
\label{tab:cer}
\begin{tabular}{lcc}
\toprule
Architecture & Val CER & Test CER \\
\midrule
TDS CNN (baseline) & -- & -- \\
CNN + LSTM          & -- & -- \\
CNN + BiLSTM        & -- & -- \\
Transformer         & -- & -- \\
\bottomrule
\end{tabular}
\end{table}

[Optional: Add a second table or figure if you varied data fraction, number of channels, or preprocessing (e.g., different transforms) and report those results here.]

% -----------------------------------------------------------------------------
\section{Discussion}
\label{sec:discussion}

[\textbf{Fill in with your observations.} Suggested points:]

\begin{itemize}
  \item \textbf{Recurrent vs.\ convolutional temporal modeling:} The TDS CNN uses a fixed receptive field; the RNN can integrate information over the full past (and future if bidirectional). We observed that [e.g., CNN+RNN improved over TDS because ... / did not improve because ...].
  \item \textbf{Bidirectional vs.\ unidirectional RNN:} Bidirectional processing uses future context, which is available in our offline setting. [Comment on whether BiLSTM helped or hurt and why.]
  \item \textbf{Transformer vs.\ RNN:} The Transformer encoder uses self-attention over the whole sequence. [Comment on CER and training stability: e.g., better long-range dependencies vs.\ higher compute or overfitting on limited data.]
  \item \textbf{Data and preprocessing:} [If you tried different data fractions, channel counts, or augmentations, summarize what helped or hurt and why.]
\end{itemize}

Limitations: We only evaluate on one subject; conclusions may not generalize across users. Training with more subjects or longer schedules might change the ranking of architectures.

% -----------------------------------------------------------------------------
\section*{References}
\label{sec:refs}

\bibliographystyle{plain}
\begin{thebibliography}{10}

\bibitem{ctc}
A.~Graves, S.~Fern\'andez, F.~Gomez, and J.~Schmidhuber.
\newblock Connectionist temporal classification: Labelling unsegmented sequence data with recurrent neural networks.
\newblock In \emph{ICML}, volume~2006, pages 369--376, 2006.

\bibitem{emg2qwerty}
V.~Sivakumar, J.~Seely, A.~Du, S.~R. Bittner, A.~Berenzweig, A.~Bolarinwa, A.~Gramfort, and M.~I. Mandel.
\newblock emg2qwerty: A large dataset with baselines for touch typing using surface electromyography.
\newblock arXiv:2410.20081, 2024.

\end{thebibliography}

\end{document}
